\section{Experiments}
\label{sec:experi}

\begin{figure*}[t]
    \setlength{\belowcaptionskip}{0pt}
	\centering
	\includegraphics[width=1.0\textwidth]{figures/method-comparison-vis-scene-0103.pdf}
	\caption{\textbf{Visual Comparisons of Scene Reconstruction and Novel-View Synthesis.}  
    Compared with per-scene optimization methods (Street Gaussians, PVG, DeformableGS, OminiRe), our ReconDrive maintains high-quality visual rendering in both scene reconstruction (Original View) and novel-view synthesis (1–3 meters of lateral movement leftward and rightward). Compared with the feed-forward method DrivingForward, ReconDrive preserves accurate geometric consistency—evident in the tree and the surroundings—and the performance improvement becomes more pronounced as the moving distance increases. Additionally, it exhibits less image distortion and blurriness, particularly in image boundary regions.
    Extended visual comparisons are provided in the Appendix.
    }\label{fig:comparison-vis-1}
\end{figure*}

This section presents a visual urban scene reconstruction benchmark comprising state-of-the-art per-scene optimization and feed-forward methods on the nuScenes~\cite{caesar2020nuscenes} dataset, and compares our approach against these baselines under multiple evaluation protocols to demonstrate its effectiveness.
\begin{table*}[htbp]
\centering
\renewcommand\arraystretch{1.0}
% \setlength\tabcolsep{3.5pt}
\resizebox{\linewidth}{!}{
\begin{tabular}{lccccccc}
\toprule
		\multirow{2}{*}{Method} &
		\multicolumn{3}{c}{Visual Scene Reconstruction} & 
		\multicolumn{3}{c}{Novel-View Synthesis} &
        Inference Speed~$\downarrow$ \\
		& PSNR~$\uparrow$ & SSIM~$\uparrow$ & LPIPS~$\downarrow$ & PSNR~$\uparrow$ & SSIM~$\uparrow$ & LPIPS~$\downarrow$ & (Per Scene) \\
\midrule
\multicolumn{8}{l}{\textit{Per-Scene Optimization Methods}} \\
Street Gaussians~\cite{yan2024street} & 29.18 & 0.8824 & 0.1658 & 22.98 & 0.6959 & 0.2948 &  31min \\
PVG~\cite{chen2023periodic} & 29.58 & 0.8839 & 0.2200 & 23.48 & 0.6919 & 0.2897 & 23min \\
DeformableGS~\cite{yang2024deformable} & 28.93 & 0.8832 & 0.1610 & 23.73 & 0.6919 & \textbf{0.2342} & 46min \\
OminiRe~\cite{chen2025omnire} & 29.42 & 0.8853 & 0.1577 & 23.01 & 0.6885 & 0.2762 & 35min \\

\hline
\multicolumn{8}{l}{\textit{Feed-Forward Methods}} \\
Drivingforward~\cite{tian2025drivingforward} & 22.83 & 0.7650 & 0.2563 & 21.88 & 0.6866 & 0.2979 & 5s \\
\textbf{ReconDrive (Ours)} & \textbf{32.66} & \textbf{0.9589} & \textbf{0.0618} & \textbf{23.99} & \textbf{0.7234} & 0.2591 & 15s \\
\bottomrule
\end{tabular}
}
\caption{\textbf{Scene Reconstruction and Novel-View Synthesis Evaluation Results on the nuScenes Dataset.}}
\label{tab:recons and nvs evaluation.}
\end{table*}

\subsection{Benchmark}
\paragraph{Dataset.}
The nuScenes~\cite{caesar2020nuscenes} is a classical autonomous driving dataset originally designed for 3D perception tasks, comprising 1,000 urban scenes. Each scene has approximately 20 s duration and contains multi-view images captured by six cameras. We use the original 700 training scenes to train the feed-forward model and select 14 representative validation scenes from the original validation split, covering diverse conditions in time (day/night), weather (sunny/rainy), driving behaviors (static/straight/turning), and traffic density. The detailed list of selected validation scenes is provided in the Appendix. Following~\cite{li2025uniscene}, we adopt a frame rate of 12 Hz as key frames. For each evaluation scene, frames at 0.5-second intervals (i.e., every 6th frame) serve as context frames for reconstruction, while the remaining frames are reserved for novel-view synthesis evaluation.

\paragraph{Evaluation Protocols and Metrics.}
We evaluate the performance of models in the following three manners. We also report the gaussian generation speed for each scene.
\begin{itemize}[leftmargin=9pt]
    \item Visual Scene Reconstruction. We use the context frames (i.e., every 6th frame) as ground truth images, and use Peak Signal-to-Noise Ratio (PSNR), Structural Similarity Index Mapping (SSIM), and Learned Perceptual Image Patch Similarity (LPIPS) as evaluation metrics.
    \item Novel-View Synthesis. We use the non-context frames (i.e., every 2nd-5th frames) as ground truth images, and also use PSNR, SSIM, and LPIPS as evaluation metrics.
    \item Novel-View Synthesis for 3D Perception. To evaluate novel-view synthesis for 3D perception, we simulate lateral camera shifts by displacing the ego-vehicle trajectory by $\pm 1$m, $\pm 2$m, and $\pm 3$m relative to the original route. We then evaluate the 3D perception performance of UniAD~\cite{hu2023planning}, which is pre-trained on the original nuScenes images at 2 Hz, using these rendered multi-view outputs across all seven offsets ($0$m and $\pm 1$m, $\pm 2$m, $\pm 3$m). For a consistent evaluation, we also provide the rendered images to UniAD at a 2 Hz frame rate. 3D perception results include 3D object detection and tracking (only vehicle type).
\end{itemize}

\paragraph{Baseline Models and Settings.}
We evaluate two categories of baseline approaches, comprising per-scene optimization methods and feed-forward methods. For the optimization-based category, we adapt the open-source DriveStudio~\cite{chen2025omnire} codebase to support the 12 Hz nuScenes dataset~\cite{caesar2020nuscenes}, as it was originally designed for the Waymo dataset~\cite{sun2020scalability}. Within this framework, we evaluate four state-of-the-art optimization methods: Street Gaussians~\cite{yan2024street}, PVG~\cite{chen2023periodic}, DeformableGS~\cite{yang2024deformable}, and OmniRe~\cite{chen2025omnire}. Each model is trained on the designated context frames for 30,000 iterations per validation scene. In the feed-forward category, we compare our work against the state-of-the-art DrivingForward~\cite{tian2025drivingforward} approach. To ensure a fair comparison under similar settings, we provide multiple frames as input to this baseline. All methods are evaluated at a resolution of $518 \times 280$ pixels, with detailed results provided in Table~\ref{tab:recons and nvs evaluation.} and Table~\ref{tab:perception-evaluation.}.

\subsection{Main Experiment Results}
\paragraph{Implementation Details.} We set the training-loss hyperparameters to $\lambda_{\text{percep}}=0.05$, $\lambda_{l2}=1.0$, $\lambda_{l1}=0.85$, $\lambda_{\text{ssim}}=0.15$, and $\lambda_{\text{scale}}=\lambda_{\text{opacity}}=0.01$. The model is optimized with AdamW using gradient accumulation (8 steps). Training proceeds in two stages: (1) Single-frame pre-training for 10 epochs (batch size 4, learning rate $2\times10^{-5}$, weight decay 0.01), where each input frame $T_i$ is also used as ground truth; (2) Dual-frame fine-tuning for 2 epochs (batch size 2). Dynamic-object masks and motion are extracted for five vehicle classes (car, truck, bus, trailer, construction vehicle). All experiments run on H800 GPUs, requiring ~8 GPU-days. Evaluation results are reported in Tab.~\ref{tab:recons and nvs evaluation.} and Tab.~\ref{tab:perception-evaluation.}.
Figure \ref{fig:comparison-vis-1} presents visual comparisons of scene reconstruction and novel-view synthesis using Scene-0103 from the nuScenes validation set. These qualitative results demonstrate ReconDrive's ability to achieve high-fidelity rendering while maintaining precise spatial alignment during view synthesis.

\paragraph{Visual Scene Reconstruction.} As shown in Tab.~\ref{tab:recons and nvs evaluation.}, per-scene optimization methods exhibit strong performance due to their iterative parameter tailoring. Specifically, PVG~\cite{chen2023periodic} and OminiRe~\cite{chen2025omnire} achieve the highest PSNR (29.58) and the highest SSIM (0.8853) respectively, while DeformableGS~\cite{yang2024deformable} lead in LPIPS. In contrast, existing feed-forward baselines~\cite{tian2025drivingforward} lag significantly (e.g., PSNR 22.83 for Drivingforward), highlighting the traditional performance gap between generalizable and scene-specific reconstruction.
ReconDrive (Ours) effectively bridges this gap. It outperforms \textit{all} baselines in pixel, structural and perceptual metrics, achieving a superior PSNR of 32.66 and SSIM of 0.9589 and LPIPS of 0.0618. 

\paragraph{Novel-View Synthesis Evaluation.} As shown in Tab.~\ref{tab:recons and nvs evaluation.}, per-scene optimization methods generally outperform prior feed-forward baselines in novel-view synthesis, as iterative parameter tuning better captures complex view-dependent variations. For instance, DeformableGS~\cite{yang2024deformable} and Street Gaussians~\cite{yan2024street} achieve PSNRs of 23.73 and 22.98, respectively, while the feed-forward Drivingforward~\cite{tian2025drivingforward} reaches only 21.88.
ReconDrive effectively bridges the performance gap between efficiency and quality, establishing a new state-of-the-art for feed-forward 4DGS. While maintaining a comparable perceptual quality (LPIPS) to the top-performing baseline, ReconDrive notably surpasses DeformableGS~\cite{yang2024deformable} by 0.26 dB in PSNR and 0.0315 in SSIM.. Furthermore, ReconDrive demonstrates substantial gains over existing feed-forward approaches, exceeding Drivingforward by 2.11 dB in PSNR and 0.0368 in SSIM, while significantly reducing perceptual error (LPIPS) by 0.0388. These results underscore that ReconDrive not only eliminates the traditional trade-off between inference speed and synthesis quality but also provides superior structural and perceptual consistency for novel-view generation in complex urban scenes.

\paragraph{3D Perception Evaluation.}
\begin{wraptable}{r}{0.52\textwidth}
\vspace{-8pt}
\centering
\renewcommand\arraystretch{1.0}

% ===================== Table 1 ====================
\vspace{-6pt}
\resizebox{1.0\linewidth}{!}{
\begin{tabular}{lcc}
\toprule
\multirow{2}{*}{Method} &
\multicolumn{1}{c}{Detection} & 		\multicolumn{1}{c}{Tracking} \\
& mAP (\%)~$\uparrow$  & AMOTA (\%)~$\uparrow$ \\
\midrule
\multicolumn{3}{l}{\textit{Per-Scene Optimization Methods}} \\
Street Gaussians~\cite{yan2024street} & 14.6 & 11.9 \\
PVG~\cite{chen2023periodic} & 18.5 & 14.4 \\
DeformableGS~\cite{yang2024deformable} & 16.4 & 13.4 \\
OminiRe~\cite{chen2025omnire} & 16.1 & 12.9 \\
\hline
\multicolumn{3}{l}{\textit{Feed-Forward Methods}} \\
Drivingforward~\cite{tian2025drivingforward} & 23.4 & 13.3 \\
\textbf{ReconDrive (Ours)} & \textbf{26.7} & \textbf{18.9} \\
\bottomrule
\end{tabular}
}
\caption{\textbf{3D Perception Results on Novel-View Rendered Scenes with 0 m, ±1 m, ±2 m, and ±3 m View Offsets.}}
\vspace{-6pt}
\label{tab:perception-evaluation.}
\end{wraptable}

Analysis of the 3D perception results in Table 2 reveals that per scene optimization methods, such as Street Gaussians~\cite{yan2024street}, exhibit limited efficacy with mAP $\le$ 14.6\% and AMOTA $\le$ 11.9\%. While feed forward approaches generally excel in detection—evidenced by DrivingForward~\cite{tian2025drivingforward} achieving an mAP of 23.4\%—they often fall short in tracking consistency, as seen in the AMOTA of 13.3\% for DrivingForward compared to 14.4\% for PVG~\cite{chen2023periodic}. In contrast, ReconDrive achieves state of the art performance across all metrics, reaching 26.7\% mAP and 18.9\% AMOTA. These results underscore ReconDrive’s robust 3D perception capabilities, demonstrating superior detection accuracy and tracking consistency even when subjected to novel view offsets.

\paragraph{Inference Efficiency.} To evaluate scalability, we calculate the average time cost for Gaussian generation per scene ($\approx$20 s duration). In ReconDrive, as adjacent temporal segments share context frames, we use caching mechanism to eliminate redundant computation during inference. With this optimization, ReconDrive achieves an inference speed of 15s per scene—orders of magnitude faster than per-scene optimization methods~\cite{yan2024street, chen2023periodic}, which require about 30 minutes. While slightly slower than existing feed-forward baselines~\cite{tian2025drivingforward} (5s), ReconDrive offers a superior balance between high-fidelity reconstruction and computational efficiency.